% !TeX root = ../main.tex
\chapter{Discussion and Conclusions}
\label{chap:discussion_conclusions}

This chapter resolves the central research questions and evaluates the formal hypotheses of the study, contextualizes our findings within the empirical literature, and outlines the methodological boundaries and future research trajectories of the work.

% ==============================================================================
% SECTION 6.1: ANSWERS TO RESEARCH QUESTIONS AND HYPOTHESIS VERDICTS
% ==============================================================================

\section{Answers to Research Questions and Hypothesis Verdicts}
\label{sec:disc_hypotheses}

The empirical findings from Chapter~\ref{chap:empirical_results} deliver structured answers to the core questions formulated at the outset of this inquiry, leading to definitive verdicts on hypotheses H1--H5:

\subsection{Lead-Time Emergence and Monotonicity}
The central operational question asks at what lead time surface weather provides actionable forecasting skill, and whether that skill expands monotonically toward harvest. Across the twelve-horizon evaluation ($H = 12 \dots 1$), the empirical evidence establishes a clear boundary: antecedent weather observed during winter dormancy and early vegetative emergence ($H = 12 \dots 6$, November through May) contains zero predictive signal ($R^2_{\text{OOS}} \le 0.00$). Statistical models trained on early-season conditions cannot anticipate summer moisture anomalies, and unconstrained specifications overfit on cold-season noise.

Actionable predictive skill emerges decisively in June ($H=5$), when formal Diebold--Mariano tests confirm that weather features reduce squared forecast loss significantly over secular trend baselines ($p = 0.0011$ for ElasticNet, $p = 0.00025$ for LSTM). Predictive accuracy expands rapidly during the reproductive flowering and pod-filling stages (July--August, $H=4 \dots 3$), where out-of-sample explained variance rises to $0.300$ for ElasticNet and $0.290$ for LSTM. At harvest closure ($H=1$), integrated ElasticNet reduces county RMSE from $6.124$ to $5.049\text{ bu/ac}$ ($5.152\text{ bu/ac}$ under weather-only forcing), while the Linear-Recurrent Ensemble reaches $R^2_{\text{OOS}} = 0.378$ ($\text{RMSE} = 4.831\text{ bu/ac}$), an overall error contraction of $21.1\%$ relative to trend projections. These results provide unequivocal support for \textbf{Hypothesis 1 (Incremental Weather Skill)}.

However, the common theoretical assumption that predictive skill increases strictly monotonically up to harvest is rejected by the data, leading to the non-support of \textbf{Hypothesis 2 (Lead-Time Monotonicity)}. Predictive accuracy exhibits a distinct plateau between late August ($H=3$) and late September ($H=2$). By late August, seed number and pod mass are biologically determined across the Corn Belt; subsequent autumn weather primarily governs harvest logistics and grain dry-down without further biological yield accumulation. In northern counties, ingesting late October precipitation actually introduces atmospheric noise that degrades forecast accuracy unless lead-time cutoffs are synchronized to local crop termination ($H^*(s)$). Forecast skill evolves as a step-function across phenological transitions, departing from a smooth monotonic progression.

\subsection{Temporal Sequence Granularity and the Model Architectures}
Our second inquiry examines whether deploying complex recurrent deep learning architectures provides demonstrable value over transparent linear baselines under strict chronological validation. The primary justification for deploying Long Short-Term Memory (LSTM) networks in agricultural forecasting is their architectural capacity to ingest high-frequency temporal sequence tensors without parameter proliferation. While linear models are constrained to monthly tabular aggregations ($P = 206$ features at $H=1$) to prevent severe multicollinearity, recurrent networks process 10-day decad sequences ($T=36$) via temporal weight sharing. Model selection on validation folds validates this design, confirming \textbf{Hypothesis 3 (Temporal Aggregation Granularity)}: 10-day decad sequences achieve lower validation error than coarse 30-day blocks, resolving acute heat spikes ($\text{EDD}_{30}$) and short-duration moisture stress during flowering that monthly averages dilute.

Does this sequence modeling translate into systemic superiority over simpler linear estimators? The empirical comparison reveals two distinct findings. The sequential LSTM achieves the highest individual out-of-sample accuracy among single models ($R^2_{\text{OOS}} = 0.370$ vs $0.320$ for ElasticNet at harvest closure). However, formal Diebold--Mariano testing ($p = 0.162$) demonstrates that its performance advantage is not statistically distinguishable from zero under serial correlation controls. During early vegetative stages ($H=5, 4$), unconstrained deep networks exhibit higher estimation variance, whereas penalized linear regression ($\ell_1/\ell_2$ shrinkage) provides critical stability in low signal-to-noise regimes \citep{James2023}. On monthly tabular features, non-linear tree ensembles (Random Forest $R^2 = 0.240$, XGBoost $R^2 = 0.300$) fail to outperform ElasticNet, displaying conservative variance attenuation and extrapolation plateaus during bumper crop years ($D_{10}$). Because complex neural networks do not uniformly dominate transparent baselines, the empirical findings corroborate \textbf{Hypothesis 4 (Model Complexity and Linear Parsimony)}. The global champion remains the Linear-Recurrent Ensemble ($R^2_{\text{OOS}} = 0.378$, $p = 0.017$ over ElasticNet), demonstrating that recurrent temporal gating and linear shrinkage offer complementary strengths.

\subsection{Extreme Shock Anticipation: County Heterogeneity versus Regional Aggregates}
The final empirical question investigates whether continuous regression models can anticipate catastrophic yield shortfalls (bottom decile $D_1$), or whether they collapse toward historical conditional means. The empirical answer depends directly on the spatial scale of evaluation.

At the macro-regional aggregate level, where localized idiosyncratic noise cancels out, the predictive framework exhibits remarkable early-warning fidelity. During the historic 2012 flash drought, systemic stress is detected by late July ($H=4$) and fully quantified by late August ($H=3$), projecting a regional yield of $44.37\text{ bu/ac}$ against an observed realization of $43.55\text{ bu/ac}$. Macro-regional forecasts correctly anticipate the sign, magnitude, and spatial footprint of systemic disasters months prior to harvest.

At the individual county scale, however, predicting acute downside extremes ($D_1$) proves far more challenging. While models correctly predict the negative direction of shocks, individual county point forecasts exhibit substantial conditional attenuation toward zero, predicting a mean departure of $-2.16\text{ bu/ac}$ against an observed realization of $-8.82\text{ bu/ac}$. As demonstrated in Appendix~\ref{chap:appendix_detrending}, this attenuation represents an intrinsic mathematical property of mean squared error minimization ($\mathbb{E}[Y \mid X]$ under ~65\% unobserved non-climatic variance), persisting identically under an oracle full-sample cubic trend. Localized soil compaction, tile drainage efficiency, and microclimatic extremes remain unobserved at county scale, causing $L_2$ estimators to shrink predictions toward the mean to avoid variance penalties. Consequently, \textbf{Hypothesis 5 (Tail-Risk Anticipation)} is \textbf{supported with qualifications}: continuous empirical models deliver robust directional early warning and high macro-regional precision, but county-level tail estimation requires dedicated extreme-value or quantile calibrations.

% ==============================================================================
% SECTION 6.2: CONFRONTATION WITH EXISTING LITERATURE
% ==============================================================================

\section{Methodological Discussion and Placement in the Literature}
\label{sec:disc_literature}

Our findings contribute to the empirical crop yield literature by establishing clear baselines, isolating physical signals from technological trends, and addressing common methodological ambiguities across four core dimensions:

\subsection{Operational Perimeter: Pure Weather Forcing Across a 75-Year Panel}
A central design decision of this study is restricting the predictor space exclusively to gridded surface reanalysis (ERA5-Land) across an unbroken 75-year panel (1951--2025) and 135 balanced Midwestern counties. This contrasts with recent studies that incorporate satellite remote sensing such as NDVI, Near-Infrared Reflectance of vegetation (NIRv), or Solar-Induced Chlorophyll Fluorescence \citep[SIF;][]{Xie2025, Yinetal2026}. While satellite proxies provide valuable direct observations of canopy greenness and photosynthetic activity, optical modalities restrict empirical evaluations to post-2000 satellite missions and suffer from missing observations during persistent cloud cover. Similarly, controlled trial studies incorporate granular agronomic management, irrigation schedules, or crop genetics \citep{Shooketal2021}, which are rarely available at commercial scale over multidecadal horizons.

Furthermore, operational viability requires strictly avoiding circularity or look-ahead bias. By design, our setup evaluates predictions using solely observed meteorological data realized up to each monthly cutoff date $T$, without attempting to inject dynamical seasonal climate forecasts (such as ECMWF SEAS5 or NOAA CFSv2) or synthetic weather completions. This choice prevents numerical forecast drift and calibration biases from contaminating the empirical evaluation. Our results thus delineate the theoretical upper bound of pure weather-driven predictability available to a decision-maker operating strictly in real time.

\subsection{Target Formulation and the Mechanics of the Trend Confound}
A prominent ambiguity in agricultural machine learning concerns target formulation. Studies that benchmark models directly on raw yield levels frequently report explained variances exceeding $R^2 \ge 0.65$--$0.75$, with deep autoregressive models attributing a major fraction of skill to prior-year yield lags \citep{ChenZhang2026}. 

Our controlled ablation clarifies this mechanism: in a multidecadal time series, raw yield $y_{t-1}$ exhibits a correlation exceeding $0.80$ with $y_t$ primarily because both series share deterministic technological momentum (advances in seed genetics, synthetic fertilizers, planting densities, and farm machinery). Including raw lagged yield leaks this secular trend into the feature space. When the target and prior-year yield are strictly detrended within training folds alone ($\epsilon_{t-1} = y_{t-1} - \hat{\tau}_{t-1}$), the genuine agronomic carryover signal (deep soil moisture memory, crop rotational benefits) contributes a modest $+2.8$ percentage points in $R^2_{\text{OOS}}$ for the ensemble (Table~\ref{tab:horizon_performance}). Our reconstructed raw yield fit ($R^2_{\text{level}} = 0.673$) demonstrates that high raw $R^2$ scores predominantly reflect deterministic technology, separating it from interannual weather sensitivity. Evaluating strictly in anomaly space aligns our framework with rigorous benchmarking standards \citep{Hoffman2020}, isolating actionable climate risk from technological drift.

\subsection{Temporal Resolution as the Source of Inductive Advantage}
In tabular monthly benchmarks, our empirical comparisons show that non-linear tree ensembles (Random Forest, XGBoost) do not systematically improve upon penalized linear regression (ElasticNet). When domain-informed feature engineering explicitly constructs non-linear physiological stress indicators ($\text{EDD}_{30}$, $VPD$, moisture balances), linear combinations effectively partition anomaly space without suffering from tree variance attenuation or step-function extrapolation limits in bumper crop years. This observation corroborates the systematic review of \citet{Sharmaetal2025}, which highlights that tree algorithms often provide marginal or null gains over well-specified regularized linear models in crop yield prediction.

The meaningful inductive advantage of deep learning emerges when sequential architectures leverage higher temporal granularity. While tabular models aggregate daily weather into 30-day monthly blocks to avoid severe parameter dilution ($P = 206$ at $H=1$), recurrent neural networks (LSTM) deploy temporal weight sharing across 10-day decad sequences ($T=36$). This sub-monthly resolution resolves acute, short-duration heat spikes and moisture deficits during flowering and pod fill that coarse monthly averages smooth out. The $+5.0$ percentage point advantage of the integrated LSTM ($R^2_{\text{OOS}} = 0.370$) over ElasticNet ($0.320$) derives from temporal sequence representation, retaining a parsimonious parameter structure.

\subsection{Spatial Aggregation: Reconciling County Heterogeneity with Regional Volume Forecasts}
Existing literature presents an apparent dichotomy between studies reporting modest county-level explained variance ($R^2 \approx 0.30$--$0.40$) and those reporting state- or national-level accuracy exceeding $R^2 \ge 0.75$ with RMSE below $3.0\text{ bu/ac}$ \citep{Khaki2020}. Our dual evaluation reconciles these findings as an inevitable consequence of observational scale.

County-level training ($N = 10{,}125$ historical evaluations) is statistically necessary to provide sufficient degrees of freedom to identify non-linear agronomic responses without sample starvation. At the county level, models must predict local microclimatic variance and unobserved field heterogeneity, yielding an out-of-sample RMSE of $4.83\text{ bu/ac}$ ($R^2_{\text{OOS}} = 0.378$). However, when individual county forecasts are aggregated to the macro-regional scale, independent local errors cancel via the law of large numbers ($\sigma / \sqrt{N}$). Across the 135 study counties, regional aggregate RMSE drops to $2.15\text{ bu/ac}$, regional explained variance rises to $R^2_{\text{macro}} = 0.72$, and linear correlation with realized regional production anomalies reaches $r = 0.910$ with LSTM ($0.880$ with Ensemble), rising to $r = 0.954$ ($0.946$ for Ensemble) in reconstructed physical yield space. Disaggregated county modeling captures spatial dispersion, while spatial aggregation delivers high-precision volumetric forecasts for commercial supply-chain operations.

% ==============================================================================
% SECTION 6.3: LIMITATIONS AND FUTURE RESEARCH TRAJECTORIES
% ==============================================================================

\section{Limitations and Future Research Trajectories}
\label{sec:disc_limitations}

The empirical setup was designed to isolate pure atmospheric predictability under strict out-of-sample discipline. Acknowledging its methodological boundaries highlights both inherent constraints and high-value research extensions:

\subsection{Scope Boundaries: Deliberate Exclusion of Optical Remote Sensing}
A primary boundary condition of this work is the omission of satellite remote sensing (such as NDVI, SIF, or thermal canopy temperature). Remote sensing indices are well-established in the recent literature and directly capture real-time vegetative development, pest outbreaks, and local soil drainage heterogeneity that surface weather observations cannot detect \citep{Xie2025, Yinetal2026}. Excluding optical observations was a deliberate design choice: multispectral satellite records are practically limited to the post-2000 era, whereas our core research objective required an unbroken 75-year historical baseline (1951--2025) to sample rare climatic shocks (such as the 1983 and 1988 droughts) and ensure robust statistical estimation across historical climate regimes. Nonetheless, for operational systems focused exclusively on contemporary campaigns, integrating high-resolution satellite reflectance alongside reanalysis features represents an essential direction to resolve within-county agronomic variance.

\subsection{Real-Time Forecasting Limits and Dynamical Climate Ensembles}
In real-time commercial applications, agribusinesses require early yield expectations months before reproductive growth occurs. Pre-sowing predictability is especially valuable for crop allocation decisions: if reliable advance intelligence indicated a catastrophic soybean season, growers could alter acreage allocations toward alternative crops or adjust input intensity. However, our continuous regression results establish that observed surface weather provides zero predictive signal prior to June ($H \ge 6, R^2_{\text{OOS}} \le 0.00$). 

To bridge this early-season information gap, operational forecasting systems require forward-looking meteorological projections. Simply substituting unobserved growing-season weather with historical climatological averages or naive linear completions provides little value and risks introducing severe estimation bias. A promising trajectory involves coupling empirical crop models with bias-corrected dynamical seasonal climate forecast ensembles, such as ECMWF SEAS5 or NOAA CFSv2 \citep{Ceglar2021}. However, as documented in climate literature, seasonal numerical forecasts suffer from substantial precipitation dispersion and regional calibration biases, which must be systematically corrected before their integration into empirical yield models can yield positive operational skill.

\subsection{Categorical Regime Classification for Early Operational Decision-Support}
While continuous point regression fails during pre-sowing and early vegetative windows, agricultural decision-makers often prioritize directional regime classifications over exact bushel quantities. A valuable methodological extension is reformulating early-season prediction as an extreme shortfall classification task \citep{Iizumietal2021, Vijverberg2023}. Testing whether probabilistic classifiers can discriminate between normal seasons and severe shortfall regimes before planting could provide actionable risk bounds for crop insurance design, forward contracting, and acreage allocation, even when continuous anomaly regression remains bounded at zero.

\subsection{Structural Trend Breaks and Real-Time Detrending Trade-Offs}
A salient empirical limitation revealed during out-of-sample testing is the presence of technological regime shifts. As documented in Appendix~\ref{chap:appendix_detrending}, Midwestern soybean yields display a structural break around 2014, with historical linear yield gains plateauing across the test partition. In our operational setup, respecting temporal causality and strictly preventing data snooping required selecting the target detrending specification based solely on pre-test historical data (1981--2010), where linear OLS minimized out-of-sample RMSE ($5.767\text{ bu/ac}$) and higher-order polynomials overfit. 

Ex-post, an oracle evaluator observing the full 1951--2025 sample would observe that a cubic trend eliminates secular test bias ($+0.06\text{ bu/ac}$ vs $-2.41\text{ bu/ac}$ under linear OLS). This discrepancy highlights a fundamental trade-off in real-time forecasting: recursive train-only linear models cannot anticipate abrupt technological decelerations in advance. While our sensitivity analysis in Appendix~\ref{chap:appendix_detrending} confirms that model rankings, lead-time dynamics, and relative shock detections remain invariant to trend specification, developing adaptive detrending estimators that detect structural breaks in real time remains a key priority for multi-decadal forecasting systems.
